\begin{proofbox}
	\textbf{\textcolor{CProof}{思路提示}}
	本题证明只需利用垂心定义即可：三条高线交于一点，从而每条高线对应的垂直关系自然成立。
	\medskip
	\textbf{\textcolor{CProof}{正式推导}}
	\begin{description}[style=nextline, leftmargin=2.8em, labelsep=0.5em, itemsep=0.55em]
		\item[\textbf{步骤一（明确定义）：}] 因为 $H$ 是 $\triangle ABC$ 的垂心，所以 $H$ 同时落在三边的高线上。
			\par\textbf{理由：} 垂心是三条高线的公共交点。
		\item[\textbf{步骤二（AH垂直BC）：}] 由于 $AH$ 所在直线是过 $A$ 的高线，故 $AH\perp BC$。
			\[
				AH\perp BC.
			\]
			\par\textbf{理由：} 高线定义为过顶点且垂直于对边。
		\item[\textbf{步骤三（BH垂直AC）：}] 同理，$BH\perp AC$。
			\[
				BH\perp AC.
			\]
			\par\textbf{理由：} 同步骤二。
		\item[\textbf{步骤四（CH垂直AB）：}] 同理，$CH\perp AB$。
			\[
				CH\perp AB.
			\]
			\par\textbf{理由：} 同步骤二。
		\item[\textbf{步骤五（向量形式）：}] 由垂直可得对应向量内积为零。
			\[
				\overrightarrow{AH}\cdot\overrightarrow{BC}=0,\quad \overrightarrow{BH}\cdot\overrightarrow{AC}=0,\quad \overrightarrow{CH}\cdot\overrightarrow{AB}=0.
			\]
			\par\textbf{理由：} 向量垂直的充要条件是点积为零。
	\end{description}
	\medskip
	\textbf{\textcolor{CProof}{结论回扣}}
	因此，垂心到各顶点的连线垂直于对边，且向量形式等价。
\end{proofbox}
